\section{Correctness Properties}
\label{sec:rs-correctness}

This section establishes three correctness properties for the Recursive Spawning Protocol: drift prevention, chain integrity, and termination. Each is stated formally, proved sketchedly, and connected to the BX3 Framework layer that enforces it. Together they constitute the RSP's overall guarantee: a spawn chain is an accountable, bounded, and auditable refinement process that terminates in human judgment.

% -------------------------------------------------------
\subsection{Drift Prevention}
\label{sec:rs-drift-prevention}

\medskip
\noindent\textbf{Theorem (Drift Prevention).} Let $C = \langle n_0, n_1, \ldots, n_k \rangle$ be an RSP chain rooted at $n_0$ with human anchor $h(n_0) = h(n_k)$. Then for all $i$ ($0 \leq i \leq k$), $R(n_i)$ is a descendant refinement of the intent of $h(n_0)$. Equivalently: no node in $C$ can exhibit behavior outside the authorized intent of its human anchor.

\medskip
\noindent\textbf{Proof sketch.} The proof proceeds by structural induction on chain depth, where each inductive step is enforced at runtime by the conjunction of three BX3-layer guarantees.

\medskip
\textit{Base case ($i=0$).} $R(n_0)$ is defined by $h(n_0)$ at provisioning. By the Purpose Layer's definition of authorized operating scope, $R(n_0)$ is a descendant refinement of $h(n_0)$'s intent by construction. No drift exists at the root.

\medskip
\textit{Inductive step.} Assume $R(n_i)$ is a descendant refinement of $h(n_0)$'s intent for some $i$ with $0 \leq i < k$. For $n_i$ to spawn $n_{i+1}$, two conditions must hold before the spawn is recorded:

\begin{enumerate}
  \item[(a)] $R(n_{i+1}) \subseteq R(n_i)$ --- the scope refinement constraint, required by the Purpose Layer's spawn authorization policy $P_{\mathrm{spawn}}$.
  \item[(b)] $\mathrm{depth}(n_{i+1}) \leq D_{\mathrm{max}}$ --- the depth constraint, enforced by the Bounds Engine and independently verified by the Fact Layer pre-commit hook.
\end{enumerate}

The Fact Layer pre-commit hook rejects any spawn violating (a) or (b). No actor, including the Bounds Engine, can execute a rejected spawn. Therefore $R(n_{i+1}) \subseteq R(n_i)$ holds by structural enforcement, not by policy convention.

By the inductive hypothesis, $R(n_i)$ is a descendant refinement of $h(n_0)$'s intent. Since $R(n_{i+1}) \subseteq R(n_i)$, transitivity of subset gives that $R(n_{i+1})$ is also a descendant refinement of $h(n_0)$'s intent. Hence $n_{i+1}$ does not drift. By induction, no node in $C$ drifts. $\square$

\medskip
The drift prevention guarantee rests on the conjunction of three structural conditions, each enforced by a different BX3 layer:

\begin{enumerate}
  \item[(1)] \textbf{Immutable parent pointer (Fact Layer).} $\alpha(n)$ is established once at provisioning and is thereafter immutable under the Fact Layer write protocol. Scope inheritance is traceable to a unique authorized lineage, not a manipulable graph. If $\alpha(n)$ were mutable, a drifted actor could rewrite its ancestry to launder authorization retroactively.

  \item[(2)] \textbf{Forensic ledger (Fact Layer).} Every scope assignment is recorded in the append-only, hash-chained ledger. The ledger makes the inductive proof empirically verifiable: each step of the refinement chain is a recorded, tamper-evident ledger entry. Tampering breaks hash continuity and is detectable by any observer with ledger read access.

  \item[(3)] \textbf{Chain terminates at Purpose Layer (Purpose Layer).} The exclusive human terminus closes the inductive argument at a finite bound. There is no machine-actor alternative terminus that could absorb exceptions and allow drift to compound silently. The human anchor $h(n)$ is non-collapsing: $h(n_i) = h(n_0)$ for all $n_i \in C$.
\end{enumerate}

The conjunction of (1), (2), and (3) is necessary and sufficient. If any condition is absent, drift prevention fails: a mutable parent pointer enables retroactive chain manipulation; an absent forensic ledger makes scope refinement unverifiable; a machine-actor terminus allows exceptions to be absorbed rather than escalated. With all three conditions present, drift is architecturally impossible because the immutable parent pointer prevents retroactive scope manipulation, the forensic ledger makes every refinement step auditable, and the Purpose Layer terminates the chain in human judgment, closing the accountability loop.

% -------------------------------------------------------
\subsection{Chain Integrity}
\label{sec:rs-chain-integrity}

\medskip
\noindent\textbf{Theorem (Chain Integrity).} For any RSP chain $C = \langle n_0, n_1, \ldots, n_k \rangle$, the parent pointer sequence $\alpha(n_i) = n_{i-1}$ holds for all $1 \leq i \leq k$, and no node configuration in $C$ has been modified after provisioning.

\medskip
\noindent\textbf{Proof sketch.} The proof has two components.

\medskip
\textit{Topological integrity.} Each spawn event $n_{i-1} \xrightarrow{\mathrm{spawn}} n_i$ establishes $\alpha(n_i) = n_{i-1}$ as an immutable Fact Layer property at provisioning time. The Fact Layer write protocol rejects any operation attempting to modify $\alpha(n_i)$ after provisioning. Since the chain is constructed exclusively through authorized spawn events and no parent pointer is modified after establishment, $\alpha(n_i) = n_{i-1}$ holds for all $i \geq 1$ by construction. There is no privileged actor --- including the Purpose Layer after provisioning --- capable of overwriting $\alpha(n)$. This eliminates the re-parenting drift pathway (Failure Mode 3) by making retroactive chain rewriting structurally impossible.

\medskip
\textit{Node immutability.} Each node $n_i$ is provisioned with a fixed configuration $(R(n_i), P_{\mathrm{spawn}}(n_i), h(n_i))$ established by the Purpose Layer and recorded in the Fact Layer authorization ledger. The Fact Layer enforces immutability of these parameters after provisioning. The write protocol rejects any modification request from any actor --- including the Bounds Engine and the Purpose Layer after provisioning. A \texttt{RESOLVE} directive redirects node behavior but is recorded as a new ledger entry (a state transition), not as an alteration of the provisioned configuration. The node's original configuration remains auditable in the ledger as the authoritative record of what was authorized at provisioning.

The forensic ledger's hash-chaining provides evidentially for chain integrity: any attempt to insert a falsified spawn event or reorder existing entries breaks hash continuity and is detectable by any observer with ledger read access. Chain integrity is a property of the entire historical record, enforceable at any future time --- not merely a property of the current chain state. $\square$

% -------------------------------------------------------
\subsection{Termination}
\label{sec:rs-termination}

\medskip
\noindent\textbf{Theorem (Termination).} Every RSP chain $C = \langle n_0, n_1, \ldots, n_k \rangle$ terminates --- either at a resolution state or at the escalation threshold --- within a finite number of spawn steps bounded by $\max(D_{\mathrm{max}}, D_{\mathrm{ESCALATE}})$.

\medskip
\noindent\textbf{Proof sketch.} The proof establishes two lemmas and combines them.

\medskip
\textit{Lemma 1 (Depth-bounded growth).} The chain grows only through spawn events. By the Bounds Engine depth enforcement and the Fact Layer pre-commit hook, any spawn producing $|C| \geq D_{\mathrm{max}}$ is rejected. Therefore the chain cannot exceed depth $D_{\mathrm{max}} - 1$ through autonomous spawn. Rejected spawns are logged; the Bounds Engine transitions to \texttt{ESCALATING}. Hence autonomous chain growth is bounded by $D_{\mathrm{max}}$.

\medskip
\textit{Lemma 2 (Escalation-triggered halt).} If the chain reaches $D_{\mathrm{ESCALATE}}$ before resolving the exception condition, the protocol enters \texttt{ESCALATING} state and suspends all autonomous spawn activity. The Bounds Engine awaits a directive from $h(n_k)$ --- the human anchor. No further spawn events occur until $h(n_k)$ issues \texttt{ACK}, \texttt{RESOLVE}, or \texttt{REJECT}. Since $h(n_k) = h(n_0)$ is a human principal, the directive arrives in finite time. The \texttt{REJECT} directive terminates the chain definitively; \texttt{RESOLVE} redirects behavior and records a new ledger entry; \texttt{ACK} permits continued operation under continued Purpose Layer authorization.

\medskip
\textit{Theorem conclusion.} Combining Lemma 1 and Lemma 2: the chain grows autonomously at most $D_{\mathrm{max}} - 1$ spawn steps. If it reaches $D_{\mathrm{ESCALATE}}$ before resolution, it halts and escalates. If it reaches $D_{\mathrm{max}}$ before resolution, the Fact Layer rejects further spawn and triggers mandatory escalation. In all execution paths, the chain reaches a terminal state within $\max(D_{\mathrm{max}}, D_{\mathrm{ESCALATE}})$ spawn steps. Since both $D_{\mathrm{max}}$ and $D_{\mathrm{ESCALATE}}$ are finite integers defined at provisioning, the chain terminates in finite time. $\square$

The termination guarantee also precludes infinite spawn-escalate cycles. When $D_{\mathrm{max}}$ is reached, the Fact Layer blocks further spawn and triggers mandatory escalation. The human anchor's \texttt{REJECT} directive terminates the chain. There is no mechanism by which the chain can circumvent this bound and re-enter autonomous spawn --- the Fact Layer pre-commit hook is structural, not configurable at runtime by any machine actor.

% -------------------------------------------------------
\subsection{Collective Guarantee}
\label{sec:rs-collective-guarantee}

Together, these three properties guarantee that a Recursive Spawning chain is an accountable, bounded, and auditable refinement process:

\begin{itemize}
  \item \textbf{Drift prevention} ensures that scope refinement is the only authorized direction of chain growth. Without it, successive scope expansions could migrate the chain's operating scope away from the human anchor's intent --- an autonomy runway within an otherwise bounded-depth chain. The depth-limit insufficiency theorem (Section~\ref{sec:drift-depth-limits}) establishes that depth bounds alone cannot prevent this.

  \item \textbf{Chain integrity} ensures that chain topology is stable and tamper-evident. Without it, an adversarial actor could manipulate chain parentage or node configuration to launder accountability or obscure the provenance of a spawned node. The immutable parent pointer and hash-chained forensic ledger make such manipulation structurally impossible.

  \item \textbf{Termination} ensures that the chain cannot grow indefinitely and cannot sustain an infinite spawn-escalate cycle. Without it, an unresolvable exception condition could drive unbounded chain growth that exceeds any human's supervisory capacity --- defeating the upstream BX3 accountability guarantee by overwhelming the human anchor.
\end{itemize}

The three properties are mutually reinforcing: drift prevention restricts the direction of chain growth, chain integrity preserves the topology established at provisioning, and termination halts the chain within a Purpose Layer-defined depth bound. The forensic ledger provides the evidentiary substrate that makes all three properties auditable at any future time. This is the operational expression of the upstream BX3 accountability guarantee applied to recursive agent population dynamics.
